\chapter*{Wprowadzenie: jak korzystać z tych notatek?}
\addcontentsline{toc}{chapter}{Wprowadzenie: jak korzystać z tych notatek?}
\setcounter{section}{0}

\begin{summarybox}
Te notatki są skryptem pomocniczym do zajęć laboratoryjnych z praktycznej
mechaniki kwantowej. Nie są pełnym podręcznikiem aksjomatycznej mechaniki
kwantowej i nie zastępują przygotowanych list zadań. Ich celem jest pomóc
studentom przejść od treści ćwiczenia laboratoryjnego do modelu, obliczeń
symbolicznych i numerycznych, wizualizacji, symulacji oraz interpretacji wyniku.
\end{summarybox}

\section{Czym te notatki są, a czym nie są?}

Te notatki są przewodnikiem po praktycznej stronie mechaniki kwantowej. Mają
wspierać zajęcia, na których student wykonuje obliczenia w Mathematice, a w
wybranych miejscach także w Pythonie. Punktem wyjścia są przygotowane listy
zadań laboratoryjnych. Notatki mają pomóc w zrozumieniu, jak takie zadania
rozwiązywać, jak kontrolować wynik i jak opisać fizykę stojącą za wykresem albo
symulacją.

Nie jest to kurs oparty na pełnym aksjomatycznym wykładzie mechaniki kwantowej.
Nie jest to również zbiór gotowych odpowiedzi do zadań. Student powinien używać
tego skryptu jako mapy: gdzie pojawia się dany model, jakie równania są
potrzebne, jakie są typowe pułapki i jak sprawdzić, czy obliczenie ma sens.

\section{Mechanika kwantowa jako teoria modeli, nie katalog wzorów}

Mechanika kwantowa w tych notatkach będzie traktowana jako język budowania
modeli. Wzory są ważne, ale nie są celem samym w sobie. Każdy wzór powinien
odpowiadać na pytanie: jaki układ opisujemy, jakie przyjmujemy założenia, jakie
wielkości liczymy i co możemy sprawdzić obliczeniowo?

Dlatego w wielu rozdziałach pojawia się ten sam schemat: najpierw sytuacja
fizyczna, potem równanie, następnie rozwiązanie analityczne albo numeryczne,
później wykres, kontrola wyniku i interpretacja. Taki schemat jest ważniejszy niż
pamięciowe opanowanie pojedynczego wzoru.

\section{Co znaczy „praktyczna” mechanika kwantowa?}

Słowo praktyczna oznacza tutaj umiejętność przejścia od sytuacji fizycznej do
równania, od równania do obliczenia, od obliczenia do wykresu, a od wykresu do
interpretacji. Student powinien umieć nie tylko uruchomić kod, ale także
rozpoznać, czy wynik ma sens fizyczny.

W praktyce oznacza to między innymi:
\begin{itemize}
    \item zapisanie Hamiltonianu dla danego modelu,
    \item określenie dziedziny i warunków brzegowych,
    \item dobranie jednostek i parametrów,
    \item wykonanie obliczeń symbolicznych albo numerycznych,
    \item przygotowanie czytelnego wykresu,
    \item sprawdzenie normalizacji, energii, transmisji albo normy,
    \item sformułowanie krótkiej interpretacji fizycznej.
\end{itemize}

\section{Rola Mathematiki i Pythona}

Mathematica jest w tym kursie traktowana jako laboratorium. Używamy jej do
rachunku symbolicznego, obliczeń numerycznych, szukania pierwiastków,
rozwiązywania równań różniczkowych oraz tworzenia czytelnych wykresów.

Python może pełnić podobną rolę, zwłaszcza w obliczeniach numerycznych i
wizualizacji. Jednak głównym narzędziem tych notatek jest Mathematica, ponieważ
pozwala wygodnie łączyć rachunek symboliczny, numerykę i prezentację wyniku w
jednym środowisku.

\begin{warningbox}
Program komputerowy nie rozumie fizyki za studenta. Mathematica może rozwiązać
równanie, ale nie powie automatycznie, czy użyto dobrych jednostek, czy warunki
brzegowe są fizyczne i czy wykres przedstawia \(\psi\), \(|\psi|\), czy
\(|\psi|^2\).
\end{warningbox}

\section{Mapa kursu}

Materiał jest ułożony tak, aby stopniowo przechodzić od prostych modeli do
bardziej złożonych symulacji:
\begin{equation}
\text{funkcja falowa}
\rightarrow
\text{studnie i bariery}
\rightarrow
\text{metody numeryczne}
\rightarrow
\text{oscylator}
\rightarrow
\text{atom wodoru}
\rightarrow
\text{komutatory i interferencja}
\rightarrow
\text{kubit}
\rightarrow
\text{dynamika czasowa}
\rightarrow
\text{ciało stałe}
\rightarrow
\text{układy 3D i projekty}.
\end{equation}

W pierwszej części uczymy się, jak normalizować funkcje falowe, rozwiązywać
proste równania własne i interpretować gęstości prawdopodobieństwa. Następnie
przechodzimy do tunelowania, macierzy przejścia i równań transcendentalnych.
Dalej pojawiają się oscylator harmoniczny, atom wodoru, moment pędu,
komutatory, przestrzeń Hilberta, kubity i układy złożone. Końcowe rozdziały
pokazują dynamikę zależną od czasu, pola magnetyczne, rachunek zaburzeń, sieci
krystaliczne, grafen i kropki kwantowe.

\section{Jak czytać równanie Schrödingera?}

Równanie Schrödingera należy czytać jako równanie modelu. Trzeba wskazać
Hamiltonian, potencjał, dziedzinę, warunki brzegowe, jednostki oraz szukane
stany. Dopiero wtedy równanie staje się zadaniem fizycznym.

Samo zapisanie równania
\begin{equation}
    \hat{H}\psi=E\psi
\end{equation}
nie wystarcza. Trzeba jeszcze wiedzieć, czym jest \(\hat{H}\), na jakiej
dziedzinie działa, jakie warunki brzegowe narzucamy i jaka wielkość ma zostać
porównana z wynikiem numerycznym albo wykresem.

\section{Jak czytać wykres funkcji falowej?}

Wykres funkcji falowej nie jest zwykłym wykresem położenia cząstki. Należy
odróżniać \(\psi\), jej część rzeczywistą i urojoną, moduł \(|\psi|\) oraz
gęstość prawdopodobieństwa \(|\psi|^2\).

Jeżeli wykres pokazuje \(|\psi(x)|^2\), powinien być interpretowany jako gęstość
prawdopodobieństwa. Jeżeli pokazuje \(\operatorname{Re}\psi(x)\), to pokazuje
część amplitudy, która może przyjmować wartości ujemne. Dobry raport musi jasno
podpisywać, co jest rysowane.

\section{Jednostki, skale i sanity checks}

W problemach nanostruktur często miesza się skale SI, elektronowolty, nanometry
i masę efektywną. Każde obliczenie powinno zawierać kontrolę jednostek oraz ocenę
rzędu wielkości wyniku.

Dla przykładu, energia w studni kwantowej zależy od \(1/L^2\). Pomylenie
nanometrów z metrami może zniszczyć wynik o czynnik \(10^{18}\). Dlatego w
notatkach często wracamy do pytania: czy wynik ma sens liczbowy?

\section{Stały schemat pracy laboratoryjnej}

\begin{summarybox}
Standardowy schemat pracy brzmi:
\[
\text{zdefiniuj model}
\rightarrow
\text{napisz równanie}
\rightarrow
\text{rozwiąż}
\rightarrow
\text{narysuj}
\rightarrow
\text{sprawdź jednostki i normę}
\rightarrow
\text{zinterpretuj}.
\]
\end{summarybox}

Ten schemat jest wspólny dla prawie wszystkich ćwiczeń. Zmienia się model i
narzędzie obliczeniowe, ale logika pracy pozostaje taka sama.

\section{Mini-przykład: jeden problem w czterech językach}

Rozważmy cząstkę w nieskończonej studni potencjału. Słownie: cząstka jest
uwięziona między dwiema idealnie nieprzenikalnymi ścianami. Matematycznie:
\begin{equation}
    V(x)=0\quad 0<x<L,
    \qquad
    \psi(0)=\psi(L)=0.
\end{equation}
Równanie własne prowadzi do energii
\begin{equation}
    E_n=\frac{n^2\pi^2\hbar^2}{2mL^2}.
\end{equation}
Komputerowo można narysować funkcje własne i gęstości prawdopodobieństwa.
Fizycznie trzeba zauważyć, że zwężenie studni zwiększa energie jak \(1/L^2\), a
stan podstawowy ma energię większą od zera.

\section{Typowy raport z ćwiczenia}

Raport powinien zawierać krótki opis teorii, definicję modelu i założeń,
równania, kod w Mathematice albo Pythonie, wykresy, kontrolę wyniku oraz
interpretację fizyczną. Sam kod bez komentarza nie jest jeszcze raportem.

Szczegółowy szablon raportu znajduje się w dodatku D. Warto traktować go jako
minimalny standard pracy laboratoryjnej.
